% Part: first-order-logic
% Chapter: introduction
% Section: first-order-logic

\documentclass[../../../include/open-logic-section]{subfiles}

\begin{document}

\olfileid{fol}{int}{fol}

\olsection{پہلے درجے دی منطق}

رسمی منطق دے پہلے تعارف توں تسی غالباً پہلے درجے
دی منطق نال واقف او۔\footnote{اصل وچ اسیں کم و بیش
ایہی فرض کردے آں!{} جے نہیں، تاں تسی ہور بنیادی
درجے دی کتاب، مثلاً \emph{forall x}
\citep{Magnus2021}، ویکھ سکدے او۔}
تسی اینوں ”مسوّراں والی منطق“ یا ”محمولی منطق“
وی آکھدے ہووگے۔ پہلے درجے دی منطق سب توں پہلے
اک رسمی زبان اے۔ یعنی ایس دی اک خاص لفظیات اے،
تے ایس دے اظہارات اوس لفظیات دیاں علامتاں دیاں
لڑیاں نیں۔ پر ہر لڑی جائز نہیں۔ جائز اظہارات
دیاں وکھ وکھ قسماں نیں: ٹرم، !!{formula}s تے
!!{sentence}s۔ ساڈی بنیادی دلچسپی پہلے درجے
دی منطق دے !!{sentence}s وچ اے: ایہ انگریزی
دے جملیاں دا رسمی ہم روپ دیندے نیں، تے انہاں
بارے اسیں اوہ سوال پچھ سکدے آں جنہاں وچ منطقی
ماہر عموماً دلچسپی رکھدا اے۔ مثال لئی:
\begin{itemize}
    \item کی $!B$،~$!A$ دا منطقی نتیجہ اے؟
    \item کی $!A$ منطقی طور تے سچا، منطقی طور تے
غلط، یا حالت اُتے منحصر اے؟
    \item کی $!A$ تے $!B$ ہم ارز نیں؟
\end{itemize}

ایہ سوال بنیادی طور تے پہلے درجے دی منطق دے
!!{sentence}s دے ”معنی“ بارے نیں۔ مثال لئی،
فلسفی ایس سوال نوں، کی $!B$ منطقی طور تے~$!A$
توں نکلدا اے، ایویں پرکھے گا: کی کوئی ایسی حالت
اے جتھے $!A$ سچا اے پر~$!B$ غلط
($!B$،~$!A$ توں نہیں نکلدا)، یا ہر اوہ حالت
جیہڑی $!A$ نوں سچا بناندی اے،~$!B$ نوں وی
سچا بناندی اے ($!B$،~$!A$ توں نکلدا اے)؟
پر ہن تک اسیں ایہ نہیں دسیا کہ ”حالت“ کیہ اے؛
ایہ کم \emph{معنویات} دا اے۔ پہلے درجے دی منطق
دیاں معنویات فلسفی دے ”حالت“ والے فطری خیال
دا ریاضیاتی طور تے ٹھیک نمونہ دیندیاں نیں،
تے نال ایہ وی طے کردیاں نیں---تے ایہ اہم اے---
کہ کسے !!a{sentence}~$!A$ دا کسے حالت
\emph{وچ سچا ہونا} کیہ اے۔ اسیں جیہڑا
ریاضیاتی طور تے ٹھیک نمونہ بناؤن لگے آں، اوہنوں
!!a{structure} آکھدے آں۔ جیہڑا تعلق
”وچ سچا“ نوں ٹھیک معنی دیندا اے، اوہنوں
\emph{پورا اترن} دا تعلق آکھدے نیں۔ سو اسیں
ایہ تعریف کراں گے کہ $!A$،
ساخت~$\Struct{M}$ وچ پورا اتردا اے
(علامتاں وچ: $\Sat{M}{!A}$)،
جتھے !!{sentence}s~$!A$ تے
!!{structure}s~$\Struct{M}$ ہون۔
ایہ ہو جان توں بعد اسیں ”توں نتیجہ نکلنا“
یا ”منطقی طور تے سچا ہونا“ ورگیاں ہور
معنوی اصطلاحاں دیاں ٹھیک تعریفاں وی دے سکاں
گے۔ انہاں تعریفاں نال ریاضیاتی طور تے
ایہ طے ہو سکے گا کہ، مثال لئی،
% PNB-INT-001: چیونٹی/کیڑے دی دلیل دے پانچ bracket slips عین نشان زدہ درست کیتے۔
$\lforall[x][(!A(x) \lif !B(x))],
\lexists[x][!A(x)] \Entails \lexists[x][!B(x)]$
درست اے یا نہیں۔ جواب یقیناً ”ہاں“ اے۔
جے تسی انگریزی دے جملیاں نوں پہلے درجے
دی منطق وچ علامتاں نال لکھنا سکھ چکے او،
تاں ایہ، مثال لئی، ”ساریاں چیونٹیاں کیڑے
نیں؛ چیونٹیاں موجود نیں؛ ایس لئی کیڑے موجود
نیں“ دی علامتی صورت اے۔ ایہ صاف طور تے
درست دلیل اے؛ ایس لئی ساڈی رسمی زبان
وچ ”توں نتیجہ نکلنا“ دا ریاضیاتی نمونہ
وی ایہی جواب دینا چاہیدا اے۔

رسمی منطق دے پہلے تعارف توں تہانوں ایہ
وی یاد ہووے گا کہ \emph{!!{derivation}s}
ہوندے نیں۔ جے تسی رسمی منطق دا ابتدائی
کورس پڑھیا اے، تاں استاد نے تہانوں ایہوجے
!!{derivation}s لبھن دی مشق کرائی ہووے گی،
شاید اک ایسا !!a{derivation} وی لبھن
لئی آکھیا ہووے جیہڑا اُتّے والا منطقی
نتیجہ وکھاندا اے۔ !!{derivation}s دین
دے کئی طریقے نیں: ممکن اے تسی ”فطری
استنباط“ یا ”سچائی دے درخت“ پڑھے ہون،
پر ہور وی بہت نیں۔ !!{derivation} نظاماں
دا مقصد اوہ اوزار دینا اے جنہاں نال اُتّے
والے منطقی سوالاں دے جواب لبھن۔ مثال لئی،
فطری استنباط دا اک !!{derivation} جیہدے
% PNB-INT-001: universal premise تے existential conclusion دے source bracket slips درست۔
مقدمے
$\lforall[x][(!A(x) \lif !B(x))]$
تے $\lexists[x][!A(x)]$ ہون تے نتیجہ
(آخری سطر) $\lexists[x][!B(x)]$ ہووے،
ایہ \emph{تصدیق کردا اے} کہ
$\lexists[x][!B(x)]$ منطقی طور تے
$\lforall[x][(!A(x) \lif !B(x))]$
تے $\lexists[x][!A(x)]$ توں نکلدا اے۔

پر ایہ کیوں؟ بظاہر !!{derivation} نظاماں
دا معنویات نال کوئی تعلق نہیں: رسمی
!!{derivation} دینا صرف علامتاں نوں
کجھ قاعدیاں مطابق رکھنا اے؛ اوس وچ
”حالتاں“ یا ”وچ سچا“ دا ذکر وی نہیں۔
!!{derivation} نظاماں تے معنویات دا
تعلق مابعد منطقی تحقیق نال قائم کرنا
پینا اے۔ لوڑ اک ریاضیاتی ثبوت دی اے،
% PNB-INT-001: آخری biconditional مثال دے source bracket slips درست۔
مثلاً ایہ کہ
$\lexists[x][!B(x)]$ دا رسمی
!!{derivation} مقدمیاں
$\lforall[x][(!A(x) \lif !B(x))]$
تے $\lexists[x][!A(x)]$ توں
اُدوں تے صرف اُدوں
ممکن اے جدوں
$\lforall[x][(!A(x) \lif !B(x))]$
تے $\lexists[x][!A(x)]$ مل کے
~$\lexists[x][!B(x)]$ دا منطقی
نتیجہ دیندے ہون۔ ایہ کم کرنے توں
پہلاں، البتہ، نحو تے معنویات دیاں
تعریفاں ٹھیک بٹھان لئی کافی باریک
کم کرنا پینا اے۔

\end{document}
